% Źródła: Eves s. 286, 298--304 (§7.4, §7.5); Greenberg 2010, s. 208.

\section{Suma kątów, defekt i pole trójkąta} % lang-pl
\section{Somma degli angoli, difetto e area del triangolo} % lang-it
\label{sekcja_suma_katow}

Do figury Saccheriego wrócimy teraz z innymi oczami: to, co on uznał za sprzeczność, okaże się twierdzeniem. % lang-pl
W tej sekcji zobaczymy, że suma kątów trójkąta jest zawsze mniejsza od $\pi$, że trójkąty o równych kątach są przystające i że pole trójkąta jest proporcjonalne do niedomiaru sumy jego kątów -- co Lambert zauważy sto lat wcześniej. % lang-pl
Torneremo ora alla figura di Saccheri con occhi diversi: ciò che lui prese per una contraddizione si rivelerà un teorema. % lang-it
In questa sezione vedremo che la somma degli angoli di un triangolo è sempre minore di $\pi$, che i triangoli con angoli uguali sono congruenti e che l'area di un triangolo è proporzionale al difetto della somma dei suoi angoli -- cosa che Lambert avrà notato un secolo prima. % lang-it

\begin{proposition}
	W geometrii hiperbolicznej kąty przy górnej podstawie czworokąta Saccheriego są równe i ostre. % lang-pl
	Odcinek łączący środki podstawy i górnej podstawy jest prostopadły do obu. % lang-pl
	Dwa czworokąty Saccheriego o równych górnych podstawach i równych kątach przy nich są przystające. % lang-pl
	Nella geometria iperbolica gli angoli alla base superiore di un quadrilatero di Saccheri sono uguali e acuti. % lang-it
	Il segmento che congiunge i punti medi della base e della base superiore è perpendicolare a entrambe. % lang-it
	Due quadrilateri di Saccheri con basi superiori uguali e angoli superiori uguali sono congruenti. % lang-it
\end{proposition}

Dowody: Eves \cite[s. 298--299]{eves1_1972}. % lang-pl
Dimostrazioni: Eves \cite[p. 298--299]{eves1_1972}. % lang-it

\begin{theorem}
	Suma kątów każdego trójkąta jest mniejsza od dwóch kątów prostych. % lang-pl
	La somma degli angoli di ogni triangolo è minore di due angoli retti. % lang-it
\end{theorem}

Dowód (z użyciem czworokąta Saccheriego zbudowanego z odcinka łączącego środki dwóch boków): Eves \cite[s. 299--300]{eves1_1972}. % lang-pl
Konsekwencje: w każdym trójkącie co najmniej dwa kąty są ostre, a suma kątów wypukłego $n$-kąta jest mniejsza od $n - 2$ kątów półpełnych; odcinek łączący środki dwóch boków trójkąta jest krótszy od połowy trzeciego boku \cite[s. 301, 304]{eves1_1972}. % lang-pl
Dimostrazione (con un quadrilatero di Saccheri costruito a partire dal segmento che congiunge i punti medi di due lati): Eves \cite[p. 299--300]{eves1_1972}. % lang-it
Conseguenze: in ogni triangolo almeno due angoli sono acuti, la somma degli angoli di un $n$-agono convesso è minore di $n - 2$ angoli piatti; il segmento che unisce i punti medi di due lati di un triangolo è minore della metà del terzo lato \cite[p. 301, 304]{eves1_1972}. % lang-it

\begin{theorem}
	Dwa trójkąty o równych kątach są przystające. % lang-pl
	Due triangoli con angoli uguali sono congruenti. % lang-it
\end{theorem}

Dowód i uwaga, że w geometrii Łobaczewskiego nie istnieją figury podobne i nieprzystające: Eves \cite[s. 300--301]{eves1_1972}. % lang-pl
Dimostrazione e osservazione che nella geometria di Lobachevskij non esistono figure simili e non congruenti: Eves \cite[p. 300--301]{eves1_1972}. % lang-it

\begin{definition}
[defekt trójkąta] % lang-pl
[difetto di un triangolo] % lang-it
\index{defekt trójkąta} % lang-pl
\index{difetto di un triangolo} % lang-it
	Niedomiar sumy kątów trójkąta względem dwóch kątów prostych: $\pi - (\alpha + \beta + \gamma)$. % lang-pl
	Il difetto della somma degli angoli di un triangolo rispetto a due angoli retti: $\pi - (\alpha + \beta + \gamma)$. % lang-it
\end{definition}

Eves \cite[s. 300]{eves1_1972}. % lang-pl
Eves \cite[p. 300]{eves1_1972}. % lang-it

Defekt jest addytywny: przy dowolnym rozcięciu trójkąta na trójkąty defekt całości jest sumą defektów części \cite[s. 301--302]{eves1_1972}. % lang-pl
Dwa trójkąty o równych defektach są równoważne przez rozkład (można je rozciąć na przystające części) \cite[s. 302--304]{eves1_1972}, więc mają to samo pole. % lang-pl
Il difetto è additivo: per ogni scomposizione di un triangolo in triangoli, il difetto del tutto è la somma dei difetti delle parti \cite[p. 301--302]{eves1_1972}. % lang-it
Due triangoli con lo stesso difetto sono equiscomponibili (si possono tagliare in parti congruenti) \cite[p. 302--304]{eves1_1972}, dunque hanno la stessa area. % lang-it

\begin{theorem}
	W geometrii Łobaczewskiego pole trójkąta jest proporcjonalne do jego defektu. % lang-pl
	Nella geometria di Lobachevskij l'area di un triangolo è proporzionale al suo difetto. % lang-it
\end{theorem}

Eves \cite[s. 304]{eves1_1972}. % lang-pl
Eves \cite[p. 304]{eves1_1972}. % lang-it

Ta zależność ma bliźniaka w geometrii sferycznej, gdzie pole trójkąta jest proporcjonalne do nadmiaru sumy kątów (zobacz sekcję \ref{sekcja_sferyczna}). % lang-pl
Pole trójkąta jest więc w geometrii hiperbolicznej ograniczone z góry, co jest zaprzeczeniem jednego ze zdań z sekcji \ref{sekcja_dobre_stare}; z ogólną teorią pola wielokątów spotkaliśmy się w rozdziale o aksjomatyce (twierdzenie \ref{twierdzenie_wallace_bolyai_gerwien}). % lang-pl
Questa relazione ha un gemello nella geometria sferica, dove l'area di un triangolo è proporzionale all'eccesso della somma degli angoli (si veda la sezione \ref{sekcja_sferyczna}). % lang-it
L'area di un triangolo è dunque limitata superiormente nella geometria iperbolica, il che nega uno degli enunciati della sezione \ref{sekcja_dobre_stare}; della teoria generale dell'area dei poligoni ci siamo occupati nel capitolo sull'assiomatica (teorema \ref{twierdzenie_wallace_bolyai_gerwien}). % lang-it

Zadania: Eves \cite[s. 301, 304]{eves1_1972}. % lang-pl
Esercizi: Eves \cite[p. 301, 304]{eves1_1972}. % lang-it

Współcześnie wynik ten ma ogólniejszą postać, zwaną twierdzeniem o jednorodności: w każdej płaszczyźnie Hilberta wszystkie trójkąty mają sumę kątów mniejszą od $\pi$, albo wszystkie równą $\pi$, albo wszystkie większą od $\pi$ (równoważnie: czwarty kąt każdego czworokąta Lamberta jest ostry, prosty albo rozwarty). % lang-pl
Według przypisu Hilberta, cytowanego przez Greenberga, udowodni to Max Dehn, a późniejsze dowody podadzą F.~Schur i Johannes Hjelmslev (zobacz sekcję \ref{sekcja_hjelmslev}); Saccheri i Lambert mieli ten wynik wcześniej, ale pierwszy z użyciem ciągłości \cite[s. 208]{greenberg_2010}. % lang-pl
W płaszczyznach typu ostrego i rozwartego trójkąty podobne są przystające, więc nie ma w nich teorii podobieństwa \cite[s. 208]{greenberg_2010}. % lang-pl
Oggi questo risultato ha una forma più generale, detta teorema di uniformità: in ogni piano di Hilbert tutti i triangoli hanno somma degli angoli minore di $\pi$, oppure tutti uguale a $\pi$, oppure tutti maggiore di $\pi$ (equivalentemente: il quarto angolo di ogni quadrilatero di Lambert è acuto, retto oppure ottuso). % lang-it
Secondo una nota di Hilbert, citata da Greenberg, lo dimostrerà Max Dehn, e dimostrazioni successive daranno F.~Schur e Johannes Hjelmslev (si veda la sezione \ref{sekcja_hjelmslev}); Saccheri e Lambert lo avevano prima, ma il primo con l'uso della continuità \cite[p. 208]{greenberg_2010}. % lang-it
Nei piani di tipo acuto e ottuso i triangoli simili sono congruenti, dunque non esiste una teoria della similitudine \cite[p. 208]{greenberg_2010}. % lang-it
